\section{Additional Application Features}
\label{sec:app_other}

\subsection{User Authentication}

Authentication is handled by a full-screen login gate rendered before the main application.
The flow proceeds in three stages: the user enters a user ID and submits; the frontend calls \texttt{POST /auth/check}; if the ID is found in MongoDB the session begins immediately, and if not, a confirmation prompt offers account creation via \texttt{POST /auth/create} or the option to re-enter a different ID.
A Continue as Guest button on the login page generates a random \texttt{guest\_XXXXXX} identifier and bypasses both checks.
No password is stored or required at any point.
The resolved user ID is persisted in \texttt{localStorage}; subsequent visits load it and bypass the gate automatically.

\subsection{Not Interested}

Each book card provides a Not Interested button.
Clicking it immediately removes the item from the current recommendation panels on the frontend and sends a \texttt{POST /interact} request with \texttt{action = "not\_interested"}.
The backend stores the item identifier in a per-user Redis set (\texttt{user:\{id\}:blocked}); subsequent calls to \texttt{GET /recommend} filter out any blocked items before enriching the candidate lists, so the dismissed book will not reappear across sessions.
No gradient step is triggered for not-interested events.

\subsection{Interaction Logging and Online Training}

The \texttt{POST /interact} endpoint accepts \texttt{user\_id}, \texttt{item\_id}, and \texttt{action} (\texttt{click} or \texttt{not\_interested}).
For click actions, the endpoint borrows a DIF-SASRec agent from the pool, loads the user's personal checkpoint, runs one sampled softmax gradient step, saves the updated checkpoint, and returns the resulting \texttt{sasrec\_loss} value.
Not-interested actions are written to a per-user Redis blocked set with no training step.
The frontend appends the returned loss value to the sparkline on response, providing sub-second visual confirmation that the model has updated.

\subsection{Streaming LLM Book Assistant}
\label{sec:llm_assistant}

Each book card exposes an Ask AI button that opens an inline popover and calls \texttt{POST /ask\_llm\_stream}.
The Qwen2.5-1.5B-Instruct model~\cite{qwen2025} generates a structured response in three sections — Genre, Plot Summary, and Pitch — grounded by a retrieved external description that is semantically re-ranked by BGE-M3 before being passed to the model.
The response streams token-by-token into the popover via \texttt{TextIteratorStreamer}; the first token dismisses the loading indicator and content builds progressively.

\subsubsection{Grounding Pipeline}

To reduce hallucination and produce content-specific descriptions rather than generic summaries from training knowledge, the assistant uses a two-stage grounding pipeline before LLM generation:

\begin{enumerate}
    \item Context retrieval: The system first queries the Google Books API using \texttt{intitle:} and \texttt{inauthor:} operators for precision, falling back to a Wikipedia search with progressively broader queries (volume-specific → series-level → title-only) if Google Books returns no description.
    Rate-limiting logic backs off for 120 seconds after a 429 response, and an \texttt{@lru\_cache(maxsize=128)} stores results so repeated queries for the same title cost 0~ms.
    \item Semantic re-ranking: The raw retrieved text is split into sentences and each sentence is scored by cosine similarity to a query embedding \texttt{``plot summary arc story characters events \{title\} volume \{vol\}''}.
    The top-5 highest-scoring sentences are selected and concatenated as the grounding context.
    This step, implemented using the same BGE-M3 encoder used throughout the system, suppresses author biography sentences and publisher marketing text in favour of plot-relevant content.
\end{enumerate}

\subsubsection{Prompt Design}

The Qwen2.5 model is invoked with a two-turn chat template.
The system prompt constrains the model to rewrite only the provided context without drawing on its own training knowledge:

\begin{quote}
\texttt{You are a book description assistant. Your ONLY job is to rewrite the provided book description into the output format. Do NOT use your own knowledge. Do NOT add plot points not in the description. Use ONLY what is written in the description below.}
\end{quote}

The user turn specifies the exact output format and mirrors the Genre field from the item's metadata, preventing the model from inferring or inventing genre labels:

\begin{quote}
\texttt{Rewrite the description above into this exact format:}\\
\texttt{**Genre:** \{genre\}}\\
\texttt{**Plot Summary:** [2--3 sentences using ONLY the description above]}\\
\texttt{**Pitch:** [1 sentence on why to read this specific volume]}\\
\texttt{Do not add anything not in the description.}
\end{quote}

Generation uses greedy decoding (\texttt{do\_sample=False}) with a 200-token budget.
Greedy decoding is preferred over sampling because the structured output format requires exact heading placement; temperature-sampled outputs occasionally reorder or omit headings.
If no grounding context is available (e.g., the item has no Google Books or Wikipedia entry), a static fallback response is emitted without invoking the model, avoiding a generation that would be entirely hallucinated.

\subsection{User Profile Endpoint}

\texttt{GET /profile?user\_id=\{id\}} returns aggregated engagement statistics: total interaction count, total search count, click-through rate, and a hydrated list of the ten most recent interactions with full item metadata.
The frontend uses this response to populate the \texttt{ProfileRadar} chart and the Activity History tab.

\subsection{Health Check}

\texttt{GET /health} reports the readiness status of each loaded component, including FAISS indices, encoders, and the AgentPool.
It is intended for infrastructure monitoring and can be queried by external health-check probes to determine when the server is ready to serve traffic.
